% =====================================================================
%  6. Gestión de incidentes   (proceso IR de ISO/IEC/IEEE 29119-2)
% =====================================================================
\section{Gestión de incidentes de prueba}
\label{sec:gestion-incidentes}

\begin{notanormativa}[Apartado nuevo respecto de la plantilla original]
La plantilla original no preveía gestión de incidentes. Sin ella, un resultado fallido no
produce ninguna consecuencia documentada: no puede comunicarse, priorizarse, seguirse ni
cerrarse. Este apartado desarrolla el proceso de \textbf{reporte de incidentes (IR1–IR2)} de
\normap{2}. El formato de registro está en el \ref{anx:incidentes}.
\end{notanormativa}

% ---------------------------------------------------------------------
\subsection{Análisis previo al reporte (IR1)}

\begin{notanormativa}[Un resultado fallido no es automáticamente un defecto del producto]
La actividad IR1 exige, antes de crear un reporte, determinar el origen de lo observado. Un
resultado inesperado puede deberse a:
\begin{itemize}
  \item un \textbf{defecto del objeto de prueba} —el caso que interesa al producto—;
  \item un \textbf{defecto del entorno} (configuración, versión desplegada, servicio caído);
  \item un \textbf{defecto de los datos} de prueba;
  \item un \textbf{error del propio caso de prueba} (resultado esperado mal derivado);
  \item un \textbf{incidente ya conocido}, que debe actualizarse en lugar de duplicarse.
\end{itemize}
Omitir esta discriminación infla el conteo de defectos y desvía el esfuerzo de corrección.
\end{notanormativa}

\begin{instruccion}
Indique quién realiza este análisis, en qué momento y qué evidencia debe acompañarlo. Registre
el origen determinado en la columna correspondiente del \ref{anx:incidentes}.
\end{instruccion}

\campo{Describa el procedimiento de análisis previo al reporte}

% ---------------------------------------------------------------------
\subsection{Severidad y prioridad}

\begin{notanormativa}[Son dimensiones independientes]
La \textbf{severidad} mide el impacto técnico o de negocio del incidente. La \textbf{prioridad}
mide la urgencia de atenderlo. Son independientes: un incidente puede ser de severidad baja y
prioridad alta (un error ortográfico en la pantalla de inicio antes de una demostración), o de
severidad alta y prioridad baja (un fallo grave en una función que no se usará en el periodo).
Colapsarlas en un solo campo impide priorizar el trabajo de corrección de forma defendible.
\end{notanormativa}

\begin{instruccion}
Defina las escalas antes de registrar el primer incidente y aplíquelas de forma consistente.
Indique también quién asigna cada valor: normalmente la severidad la propone quien prueba y la
prioridad la decide quien decide sobre el producto.
\end{instruccion}

\begin{table}[H]
\centering
\caption{Escala de severidad acordada.}
\label{tab:severidad}
\begin{tabularx}{\textwidth}{|C{2.2cm}|Y|C{2.4cm}|}
  \hline
  \filaenc \textbf{Nivel} & \textbf{Significado: impacto observable} & \textbf{Quién la asigna} \\ \hline
  \campo{Crítica} & \campo{Criterio} & \campo{Rol} \\ \hline
  \campo{Alta}    & \campo{Criterio} & \campo{Rol} \\ \hline
  \campo{Media}   & \campo{Criterio} & \campo{Rol} \\ \hline
  \campo{Baja}    & \campo{Criterio} & \campo{Rol} \\ \hline
\end{tabularx}
\end{table}

\begin{table}[H]
\centering
\caption{Escala de prioridad acordada.}
\label{tab:prioridad}
\begin{tabularx}{\textwidth}{|C{2.2cm}|Y|C{2.4cm}|}
  \hline
  \filaenc \textbf{Nivel} & \textbf{Significado: urgencia de atención} & \textbf{Quién la asigna} \\ \hline
  \campo{Inmediata} & \campo{Criterio y plazo} & \campo{Rol} \\ \hline
  \campo{Alta}      & \campo{Criterio y plazo} & \campo{Rol} \\ \hline
  \campo{Normal}    & \campo{Criterio y plazo} & \campo{Rol} \\ \hline
  \campo{Baja}      & \campo{Criterio y plazo} & \campo{Rol} \\ \hline
\end{tabularx}
\end{table}

Los Cuadros~\ref{tab:severidad} y~\ref{tab:prioridad} definen dos ejes distintos de
clasificación; ambos se registran, por separado, en cada incidente del \ref{anx:incidentes}.

% ---------------------------------------------------------------------
\subsection{Estados y flujo de tratamiento}

\begin{instruccion}
Defina los estados por los que puede pasar un incidente y quién autoriza cada transición. Un
estado sin responsable es un incidente que nadie mueve. Incluya explícitamente el estado
terminal que corresponda cuando el defecto \emph{no} se corrige, para que los incidentes
abiertos al cierre queden documentados en lugar de desaparecer.
\end{instruccion}

\begin{table}[H]
\centering
\caption{Estados de un incidente y responsables de cada transición.}
\label{tab:estados-incidente}
\begin{tabularx}{\textwidth}{|C{2.4cm}|Y|C{2.4cm}|C{2.4cm}|}
  \hline
  \filaenc \textbf{Estado} & \textbf{Significado} & \textbf{Transiciona a} & \textbf{Responsable} \\ \hline
  Nuevo        & Registrado, aún no analizado por el equipo receptor. & \campo{Abierto / Rechazado / Duplicado} & \campo{Rol} \\ \hline
  Abierto      & Aceptado como incidente válido y pendiente de tratamiento. & \campo{En corrección / Diferido} & \campo{Rol} \\ \hline
  En corrección& El equipo desarrollador trabaja en la resolución. & \campo{Resuelto} & \campo{Rol} \\ \hline
  Resuelto     & Se entregó una corrección; pendiente de confirmación. & \campo{Cerrado / Reabierto} & \campo{Rol} \\ \hline
  Cerrado      & Prueba de confirmación aprobada (\ref{sec:confirmacion-regresion}). & --- & \campo{Rol} \\ \hline
  Diferido     & Válido, pero no se corregirá en este periodo. & \campo{Riesgo residual} & \campo{Rol} \\ \hline
  Rechazado    & No constituye un defecto del objeto de prueba (véase IR1). & --- & \campo{Rol} \\ \hline
  Duplicado    & Ya registrado; se actualiza el incidente original. & --- & \campo{Rol} \\ \hline
\end{tabularx}
\end{table}

El Cuadro~\ref{tab:estados-incidente} define el ciclo de vida del incidente. Los estados
\emph{Diferido} y \emph{Abierto} al cierre del esfuerzo alimentan directamente los incidentes
abiertos (apartado~\ref{sec:incidentes-abiertos}) y los riesgos residuales
(apartado~\ref{sec:residuales}).

\begin{politicaacademica}
Dado que el curso establece que el equipo de prueba no repara los defectos encontrados, es
previsible que muchos incidentes terminen el periodo en estado \emph{Abierto} o
\emph{Diferido}. Eso no es un fallo del proceso: lo que sí debe evitarse es cerrarlos sin
corrección ni confirmación, porque haría irreconocible la diferencia entre un defecto resuelto
y uno simplemente no atendido.
\end{politicaacademica}

% ---------------------------------------------------------------------
\subsection{Comunicación de incidentes}

\begin{instruccion}
Indique cómo se comunican los incidentes al equipo desarrollador y a los demás interesados: por
qué medio, con qué plazo según la severidad, y quién es el punto de contacto de cada parte.
\end{instruccion}

\begin{table}[H]
\centering
\caption{Comunicación de incidentes según severidad.}
\label{tab:comunicacion-incidentes}
\begin{tabularx}{\textwidth}{|C{2.2cm}|C{3.0cm}|Y|C{2.4cm}|}
  \hline
  \filaenc \textbf{Severidad} & \textbf{Destinatario} & \textbf{Medio} & \textbf{Plazo} \\ \hline
  \campo{Crítica} & \campo{Rol} & \campo{Medio} & \campo{Plazo} \\ \hline
  \campo{Otras}   & \campo{Rol} & \campo{Medio} & \campo{Plazo} \\ \hline
\end{tabularx}
\end{table}

El Cuadro~\ref{tab:comunicacion-incidentes} fija los plazos de aviso; su cumplimiento es uno de
los aspectos que el seguimiento (apartado~\ref{sec:seguimiento}) puede monitorear.
